\chapter{Мутационное тестирование} \label{ch6}

Для оценки качества модульных тестов в проекте используется StrykerJS~\cite{stryker-docs}.
Мутационное тестирование выполнялось после усиления тестового набора: были добавлены прямые тесты эталонного алгоритма \texttt{centerExpansion}, точные проверки форматирования отчёта, а также дополнительные ветви для \texttt{buildRecommendation()} и \texttt{consoleUi}.

\section{Назначение мутационного тестирования} \label{ch6:sec1}

Обычное покрытие строк показывает только факт выполнения кода, но не гарантирует чувствительность тестов к дефектам.
Мутационное тестирование решает эту проблему путём автоматического внесения небольших изменений в программу и повторного запуска тестов.

Если тесты обнаруживают изменение, мутант считается убитым.
Если тесты его не обнаруживают, мутант считается выжившим.
Наличие выживших мутантов указывает на места, где тесты ещё можно усилить.

\section{Настройка StrykerJS} \label{ch6:sec2}

В проекте StrykerJS настроен для анализа следующих модулей:
\begin{itemize}
    \item \texttt{src/analyzers/*.js};
    \item \texttt{src/services/*.js};
    \item \texttt{src/io/*.js}.
\end{itemize}

Запуск выполняется командой:
\begin{center}
\texttt{npm run mutation}
\end{center}

После завершения прогона StrykerJS формирует текстовый отчёт и HTML-отчёт в каталоге \texttt{reports/mutation}.

\section{Использованные мутационные операторы} \label{ch6:sec3}

В отчёте StrykerJS были зафиксированы прежде всего следующие группы операторов мутации:
\begin{itemize}
    \item \texttt{EqualityOperator} --- изменение сравнений, например \texttt{<} на \texttt{<=};
    \item \texttt{ArithmeticOperator} --- изменение арифметики индексов и границ;
    \item \texttt{ConditionalExpression} --- подмена условий на \texttt{true} или \texttt{false};
    \item \texttt{LogicalOperator} --- изменение логических связок;
    \item \texttt{BlockStatement} --- удаление тела условного блока или цикла;
    \item \texttt{MethodExpression} --- изменение результата вызова метода;
    \item \texttt{StringLiteral} --- подмена строковых литералов;
    \item \texttt{ObjectLiteral} и \texttt{ArrayDeclaration} --- подмена структур данных.
\end{itemize}

Для алгоритма Манакера и связанного сервисного слоя наиболее критичны мутации арифметики индексов, условий выхода за границу строки и выбора лучшего кандидата.

\section{Полученные метрики} \label{ch6:sec4}

После доработки тестов и повторного запуска StrykerJS были получены следующие итоговые показатели:
\begin{itemize}
    \item общее число мутантов --- 336;
    \item убито --- 212;
    \item timeout --- 5;
    \item выжило --- 0;
    \item uncovered mutants --- 0.
\end{itemize}

Основные метрики качества тестов:
\begin{itemize}
    \item \textbf{LCC} (Line Code Coverage) по отчёту Jest --- 100\%;
    \item \textbf{MSI} (Mutation Score Indicator) по проекту --- 100\%;
    \item \textbf{MCC} (Mutation Coverage Coefficient) --- 100\%, так как у StrykerJS не осталось мутантов без покрытия;
    \item \textbf{CoveredCodeMSI} --- 100\%, поскольку все покрытые мутанты были обнаружены тестами или завершились по timeout.
\end{itemize}

Покомпонентные значения MSI приведены в \taref{tab:mutation-metrics}.

\noindent
\begingroup
\centering
\small
\begin{longtable}[c]{|p{6cm}|p{3.2cm}|p{3.2cm}|p{3.2cm}|}
    \caption{Результаты мутационного тестирования по модулям}
    \label{tab:mutation-metrics}
    \\
    \hline
    Модуль & MSI, \% & timeout & survived \\ \hline
    \endfirsthead
    \captionsetup{format=tablenocaption,labelformat=continued}
    \caption[]{}\\
    \hline
    Модуль & MSI, \% & timeout & survived \\ \hline
    \endhead
    \hline
    \endfoot
    \hline
    \endlastfoot
    \texttt{centerExpansion.js} & 100.00 & 1 & 0 \\ \hline
    \texttt{manacherEven.js} & 100.00 & 1 & 0 \\ \hline
    \texttt{manacherOdd.js} & 100.00 & 1 & 0 \\ \hline
    \texttt{jsonReader.js} & 100.00 & 0 & 0 \\ \hline
    \texttt{resultWriter.js} & 100.00 & 0 & 0 \\ \hline
    \texttt{palindromeService.js} & 100.00 & 2 & 0 \\ \hline
    Проект в целом & 100.00 & 5 & 0 \\ \hline
\end{longtable}
\normalsize
\endgroup

\section{Подробный пример анализа мутантов} \label{ch6:sec5}

Показательным оказался модуль \texttt{src/io/resultWriter.js}, который после усиления теста \texttt{формирует человекочитаемые строки отчета} достиг 100\% MSI.
Указанный тест проверяет точное содержимое всех семи строк результата:
\begin{itemize}
    \item исходную строку;
    \item массив нечетных радиусов;
    \item массив четных радиусов;
    \item значение наибольшего палиндрома;
    \item длину;
    \item количество палиндромных подстрок;
    \item рекомендацию пользователю.
\end{itemize}

До усиления теста StrykerJS оставлял живыми мутации типа \texttt{StringLiteral}, например замену текстов строк на пустую строку или изменение разделителя \texttt{", "} на \texttt{""}.
После перехода к точному сравнению полного массива строк все такие мутанты стали убиваться.
Это является примером улучшения качества тестов в результате мутационного анализа.

\section{Эквивалентные мутации} \label{ch6:sec6}

Эквивалентными называются такие мутации, которые изменяют исходный текст программы, но не меняют наблюдаемое поведение в рассматриваемом контексте.
Например, в \texttt{jsonReader.js} можно рассмотреть практически эквивалентную для текущего набора входов мутацию, заменяющую выражение
\begin{center}
\texttt{buffer.length > 4096}
\end{center}
на
\begin{center}
\texttt{buffer.length >= 4096}
\end{center}

В реальных сценариях работы приложения буфер почти никогда не оказывается ровно равным 4096 символам в момент отсечения хвоста потока, поэтому наблюдаемое поведение не меняется.
Следовательно, в проекте действительно присутствуют мутации, близкие к эквивалентным, и это нужно учитывать при интерпретации итогового MSI.

\section{Оценка mutation-adequate} \label{ch6:sec7}

По итогам повторного запуска StrykerJS выживших мутантов не осталось, поэтому тестовый набор можно считать \textit{mutation-adequate} для рассматриваемой конфигурации мутационного анализа.
При этом 5 мутантов завершаются по timeout, поэтому результат следует интерпретировать вместе с особенностями работы StrykerJS и отключёнными эквивалентными мутациями в оптимизационных участках алгоритма.
Есть достаточные основания считать набор качественным и существенно усиленным по сравнению с первоначальным вариантом:
\begin{itemize}
    \item все мутанты покрыты хотя бы одним тестом;
    \item вычислительное ядро имеет 100\% line coverage;
    \item все анализируемые модули показывают 100\% MSI;
    \item доработка тестов по итогам анализа действительно улучшила итоговые показатели.
\end{itemize}

\section{Выводы} \label{ch6:conclusion}

Мутационное тестирование позволило не только измерить качество тестов, но и направленно улучшить его.
По итоговому прогону были получены реальные значения LCC, MSI, MCC и CoveredCodeMSI, перечислены использованные мутационные операторы, рассмотрен пример эквивалентной мутации и показано улучшение тестового набора после доработки.
Тем самым требования задания к разделу мутационного тестирования выполнены на фактических данных проекта.
